\documentclass[10pt,a4paper]{amsart}
\usepackage[margin=19mm]{geometry}
\usepackage{amsmath,amssymb,mathtools}
\usepackage[scr=rsfs,cal=euler]{mathalfa}
\usepackage{xfrac}
\usepackage[unicode,hidelinks,bookmarks=false]{hyperref}
\newcommand{\ie}{i.e.\ }
\newcommand{\cf}{cf.\ }
\newcommand{\resp}{respectively\ }
\newcommand{\Subsection}{\subsection}
\providecommand{\nobreakdash}{\nobreak\hbox{-}}
\setlength{\emergencystretch}{2em}
\multlinegap0pt
\usepackage{amssymb}
\usepackage{mathtools}
\newtheorem{lemma}{Lemma}
\def\F{\mathbb{F}}
\title{Splitting sums of binary polynomials}
\author{Luis H. Gallardo}
\date{Source: arXiv:2602.13111v1 (CC BY 4.0)}
\begin{document}
\maketitle

\noindent\emph{Source and licence.} Splitting sums of binary polynomials. Version arXiv:2602.13111v1 (CC BY 4.0), distributed by arXiv under the Creative Commons Attribution 4.0 International licence, http://creativecommons.org/licenses/by/4.0/, as recorded in the arXiv licence metadata for this exact version and shown on the abstract page https://arxiv.org/abs/2602.13111v1. Full licence notice: \texttt{licenses/arxiv-2602.13111-CC-BY-4.0.txt}.\par\smallskip

\noindent\textit{Editorial scope.} Counted unit: the article's lemma ACEsept (source0.tex lines 304--319). The article's complete original proof of it is delivered with it (source0.tex lines 321--423, one \texttt{proof} environment of 103 lines). No mathematics is written, repaired or re-derived: the statement and the proof are the article's own lines, in the article's own notation, and no retained line is substituted or re-flowed.  Retained uncounted as the source-local context the proof needs: lines 206--215.  Nothing was omitted from any retained span, so the package carries no omission record.  No retained line contains a citation, so the wrapper carries no bibliography and no bibliography entry.  No retained line contains a figure environment, an image-inclusion command or drawing code, so no image is delivered and the package declares no asset dependency.  Editorial formatting only, in addition: an AMS \texttt{amsart} preamble that restates the article's own declarations and macros used by the retained lines, namely the environments \texttt{lemma} on separate counters, together with the standard packages those lines need.  The retained environments print in this document as Lemma 1 in order of appearance, whereas the article prints the same items with the numbering of its own full document.  The complete native source of the article as distributed in the arXiv e-print for this exact version is bundled unchanged as \texttt{sources/194616/source0.tex}.
\begin{lemma}
\label{ABC}
The only non-negative integers \( A, B, C \) satisfying
\begin{equation}
\label{xabc}
(x+1)^A + x^B = x^C
\end{equation}
in \( \F_2[x] \)
are either \( A = 2^s, B = 0, C = 2^s \) or \( A = 2^s, B = 2^s, C = 0 \), for some non-negative integer \( s \).
\end{lemma}

\begin{lemma}
\label{ACEsept}
Let \( A, B, C, D, E, F \) be non-negative integers such that
\begin{equation}
\label{xabcdef1}
x^A(x+1)^B + x^C(x+1)^D = x^E(x+1)^F
\end{equation}
in \( \F_2[x] \), with \( A \leq C \leq E \). If \( C = A \), then \( B \geq D \).

Then \( \min(B, D, F) = D \). Moreover,
   
     \( B \neq D \), \( C = A \), \( F = D \), and:
    
         \( B - D = E - A = 2^s \) for some non-negative integer \( s \).
       
\end{lemma}

\begin{proof}
Assume that \eqref{xabcdef1} holds with \( A \leq C \leq E \), and if \( C = A \), then \( B \geq D \). Let us write
\[
P := x^A(x+1)^B, \quad Q := x^C(x+1)^D, \quad R := x^E(x+1)^F.
\]

We proceed by dividing the equation
\[
P + Q = R
\]
by the appropriate power of \( x+1 \) to reduce the minimal exponent among \( B, D, F \) to zero.

\textbf{Step 1: Normalize by minimal power of \( x+1 \).}

Let \( m = \min(B, D, F) \). Dividing both sides of the equation by \( (x+1)^m \), we obtain
\begin{equation}
\label{eq:reduced}
x^A(x+1)^{B - m} + x^C(x+1)^{D - m} = x^E(x+1)^{F - m}.
\end{equation}
We claim that \( m = D \). Suppose not.

\textbf{Case 1: \( m = B < D \).} Then in \eqref{eq:reduced}, the left-hand side becomes
\[
x^A + x^C(x+1)^{D - B}.
\]
Since \( D - B \geq 1 \), this is a sum of two split polynomials with distinct powers of \( x+1 \). But the right-hand side \( x^E(x+1)^{F - B} \) is also a split polynomial. This contradicts Lemma \ref{ABC}, which classifies when such a sum equals a single split polynomial. Hence, this case cannot occur.

\textbf{Case 2: \( m = F < D \).} Then the right-hand side becomes \( x^E \), a monomial. But then the left-hand side is a sum of two split polynomials, which again cannot equal a monomial unless one of them is zero, which is excluded by assumption. Thus, this case also leads to a contradiction.

Therefore, we must have \( m = D \), as claimed.

\textbf{Step 2: Analyze according to whether \( B = D \) or \( B \neq D \).}

We now consider \eqref{xabcdef1}
with \( D \) minimal among \( B, D, F \). Divide both sides by \( (x+1)^D \), yielding
\begin{equation}
\label{eq:step2}
x^A(x+1)^{B - D} + x^C = x^E(x+1)^{F - D}.
\end{equation}

\textbf{Subcase 1: \( B = D \).}

Then \eqref{eq:step2} becomes
\begin{equation}
\label{eq:step2A}
x^A + x^C = x^E(x+1)^{F - D}
\end{equation}
so that $A < C \leq E$. Taking degrees into  \eqref{eq:step2A} we obtain that
\[
C = E + (F - D) \geq E.
\]
Thus, $C=E$, and $F=D$. Hence, \eqref{eq:step2A} gives the contradiction 
\[
x^A+x^C = x^C.
\]
In other words, the case $B=D$ does not happen.


\textbf{Subcase 2: \( B > D \).}

Then \eqref{eq:step2} becomes
\[
x^A(x+1)^{B - D} + x^C = x^{E}(x+1)^{F - D}.
\]

 Now two cases arise:

\begin{itemize}
    \item If \( C = A \), then \eqref{eq:step2} becomes
    \begin{equation}
    \label{eq:beefe}
    (x+1)^{B - D} + 1 = x^{E - A}(x+1)^{F - D}.
    \end{equation}
    
    We claim that $B - D > F -D \geq 0$. Otherwise, \eqref{eq:beefe} gives the contradiction
  \[
    (x+1)^{B - D} \mid 1.
    \]
    
    Thus, \eqref{eq:beefe} implies that $(x+1)^{F - D} \mid 1$, so that $F = D$.
    Therefore,  \eqref{eq:beefe} becomes
    \begin{equation}
    \label{eq:paraL6}
     1 = x^{E - A} + (x+1)^{B - D}.
    \end{equation}
    
     
       
   Apply Lemma \ref{ABC} to  \eqref{eq:paraL6}, which implies  \( B - D = E - A = 2^s \) for some \( s \geq 0 \).
   

   \item If \( C \neq A \), then we can write \eqref{eq:step2} as follows:
   
   \begin{equation}
    \label{eq:aete}
     x^A \left( (x+1)^{B-D} + x ^{C-A} \right) = x^{E} (x+1)^{F - D}.
    \end{equation}
    
   Comparing the exponent of $x$ in both sides of \eqref{eq:aete} we reach the contradiction $A = E$. In other words, the case $C \neq A$ does not happen.
 
\end{itemize}
This completes the analysis of all cases, and hence the proof.
\end{proof}

\end{document}
